\documentclass[10pt,a4paper]{amsart}
\usepackage[margin=19mm]{geometry}
\usepackage{amsmath,amssymb,mathtools}
\usepackage[scr=rsfs,cal=euler]{mathalfa}
\usepackage{xfrac}
\usepackage[unicode,hidelinks,bookmarks=false]{hyperref}
\newcommand{\ie}{i.e.\ }
\newcommand{\cf}{cf.\ }
\newcommand{\resp}{respectively\ }
\newcommand{\Subsection}{\subsection}
\providecommand{\nobreakdash}{\nobreak\hbox{-}}
\setlength{\emergencystretch}{2em}
\multlinegap0pt
\usepackage{dsfont}
\usepackage{braket}
\usepackage{amssymb}
\usepackage{tikz}
\usepackage{xcolor}
\usepackage{mathtools}
\usepackage{bm}
\newtheorem{proposition}{Proposition}
\newcommand{\Iso}{\mathrm{Iso}}
\newcommand{\bnd}{\mathrm{bnd}}
\newcommand{\diam}{\mathrm{diam}}
\newcommand{\loc}{\mathrm{loc}}
\newcommand{\phys}{\mathrm{phys}}
\title{Pathology-Free Real-Space Renormalization Group Theory on an Inverse Limit Space}
\author{Fabio Arz}
\date{Source: arXiv:2609.18356v1 (CC BY 4.0)}
\begin{document}
\maketitle

\noindent\emph{Source and licence.} Pathology-Free Real-Space Renormalization Group Theory on an Inverse Limit Space. Version arXiv:2609.18356v1 (CC BY 4.0), distributed by arXiv under the Creative Commons Attribution 4.0 International licence, http://creativecommons.org/licenses/by/4.0/, as recorded in the arXiv licence metadata for this exact version and shown on the abstract page https://arxiv.org/abs/2609.18356v1. Full licence notice: \texttt{licenses/arxiv-2609.18356-CC-BY-4.0.txt}.\par\smallskip

\noindent\textit{Editorial scope.} Counted unit: the article's proposition Eq1 (source0.tex lines 770--776). The article's complete original proof of it is delivered with it (source0.tex lines 778--831, one \texttt{proof} environment of 54 lines). No mathematics is written, repaired or re-derived: the statement and the proof are the article's own lines, in the article's own notation, and no retained line is substituted or re-flowed.  No further source-local context is needed: every cross-reference of every retained line resolves inside the two delivered spans.  Nothing was omitted from any retained span, so the package carries no omission record.  No retained line contains a citation, so the wrapper carries no bibliography and no bibliography entry.  No retained line contains a figure environment, an image-inclusion command or drawing code, so no image is delivered and the package declares no asset dependency.  Editorial formatting only, in addition: an AMS \texttt{amsart} preamble that restates the article's own declarations and macros used by the retained lines, namely the environments \texttt{proposition} on separate counters, together with the standard packages those lines need.  The retained environments print in this document as Proposition 1 in order of appearance, whereas the article prints the same items with the numbering of its own full document.  The complete native source of the article as distributed in the arXiv e-print for this exact version is bundled unchanged as \texttt{sources/210692/source0.tex}.
\begin{proposition}
    Let $\Phi$ be an absolutely summable interaction and $\Iso(\Lambda_\infty)\times\mathds{Z}^2$-invariant interaction on $\Lambda_\infty$. Then there is a canonical way to construct a corresponding $H^\Phi\in \mathcal{V}_\infty$. In case that $\Phi$ satisfies the stricter inequality
    \begin{equation}\label{Eq1}
        \sum_{\substack{\Gamma\subset\Lambda_\infty,|\Gamma|<\infty\\0\in\Gamma}}\diam(\Gamma)\cdot\|\Phi_\Gamma\|_\infty<\infty,
    \end{equation}
    we find that $H^\Phi\in \mathcal{V}_\phys$.
\end{proposition}

\begin{proof}
    Let us consider the coset $S_n=\{\Gamma\subset\Lambda_\infty\,:\,|\Gamma|<\infty\}/\mathcal{D}(\pi_{\infty n})$. We then define
    \begin{equation}
        H^{\Phi}_n(\sigma):=\sum_{[\Gamma]\in S_n}\Phi_\Gamma(\pi_{\infty n}^*\sigma).
    \end{equation}
    First, we quickly note that each map is well-defined as 
    \begin{equation}
        \sum_{[\Gamma]\in S_n}|\Phi_\Gamma(\pi_{\infty n}^*\sigma)|\leq\|\Phi\|_{\mathcal{B}^1}\cdot|\Lambda_n|<\infty
    \end{equation}
    and thus the series $\sum_{[\Gamma]\in S_n}\Phi_\Gamma(\pi_{\infty n}^*\sigma)$ converges. The next step is to prove that $H^\Phi\in \mathcal{V}_\infty$, i.e.\ that $\omega_{nm}H_n^\Phi=H_m^\Phi$. Let us therefore consider
    \begin{equation}
        H_n^\Phi(\pi_{nm}^*\sigma)=\sum_{[\Gamma]\in S_n}\Phi_\Gamma(\pi_{\infty m}^*\sigma).
    \end{equation}
    Here, we note that for each equivalence class $[\Gamma]\in S_m$ there exist equivalence classes $[\Gamma_1],...,[\Gamma_{b^{n-m}}]\in S_n$ such that
    \begin{equation}
        [\Gamma]=\bigcup_{i=1}^{b^{n-m}}[\Gamma_i].
    \end{equation}
    Further, we find that $(\pi_{\infty m}^*\sigma)_{\Gamma_i}=(\pi_{\infty m}^*\sigma)_\Gamma$. Thus, we conclude that
    \begin{equation}
        H_n^\Phi(\pi_{nm}^*\sigma)=b^{n-m}\sum_{[\Gamma]\in S_m}\Phi_\Gamma(\pi_{\infty m}^*\sigma)=b^{n-m}H_{m}^\Phi(\sigma)\implies\omega_{nm}H_n^\Phi=H_m^\Phi\,.
    \end{equation}
    For the second statement, let us assume that $\Phi$ satisfies (\ref{Eq1}). We then must prove that $\Delta H_n^\Phi$ satisfies the necessary bound such that $H^\Phi\in \mathcal{V}_\phys$. Let us start by considering
    \begin{equation}
        H_{n}^\Phi(\sigma)-\sum_{i=1}^bH_{n-1}^\Phi(\sigma_{Q_i})=\sum_{[\Gamma]\in S_n}\Phi_\Gamma(\pi_{\infty n}^*\sigma)-\sum_{i=1}^b\sum_{[\Delta]\in S_{n-1}}\Phi_\Delta(\pi_{\infty n-1}^*\sigma_{Q_i}).
    \end{equation}
    Let us now decompose $S_n=S^\loc_n\cup S_n^\bnd$ where
    \begin{equation}
        S_n^\loc:=\{[\Gamma]\in S_n\,:\,\pi_{\infty n}(\Gamma)\in Q_i,i\in\{1,...,b\}\},\qquad S_n^\bnd=S_n\backslash S_n^\loc.
    \end{equation}
    For any $[\Gamma]\in S_n^\loc$ there exists a corresponding $[\Delta(\Gamma)]\in S_{n-1}$ satisfying $\Phi_\Gamma(\pi_{\infty n}^*\sigma)=\Phi_{\Delta(\Gamma)}(\pi_{\infty n-1}^*\sigma_{Q_i})$. Therefore, we conclude that
    \begin{equation}
        \|\Delta _nH_n^\Phi\|_\infty\leq\sum_{[\Gamma]\in S_n^\bnd}\|\Phi_\Gamma\|_\infty=b\sum_{\substack{[\Gamma]\in S_n^\bnd\\\Gamma\cap \pi_{\infty n}^{-1}(Q_1)\ne\emptyset}}\|\Phi_\Gamma\|_\infty\,.
    \end{equation}
    For any $[\Gamma]\in S_n^\bnd,\Gamma\cap\pi_{\infty n}^{-1}(Q_1)\neq\emptyset$ there must exist at least two vertices $x\in Q_1,y\in \Lambda_n\backslash Q_1$ such that $\pi_{\infty n}^{-1}(x)\cap\Gamma\neq\emptyset,\pi_{\infty n}^{-1}(y)\neq\emptyset$. Consequently, we find that $\diam(\Gamma)\geq d(x,y)> d(x,\partial Q_1)$. Therefore, we may bound the previous sum via (by translation invariance we may choose $\tilde x$ to be any point in $\pi_{\infty n}^{-1}(x)$) 
    \begin{equation}
        \sum_{\substack{[\Gamma]\in S_n^\bnd\\\Gamma\cap\pi_{\infty n}^{-1}(Q_1)\ne\emptyset}}\|\Phi_\Gamma\|_\infty\leq\sum_{x\in Q_1}\sum_{\substack{\Gamma\ni\tilde x\\\diam(\Gamma)> d(x,\partial Q_i)}}\|\Phi_\Gamma\|_\infty=\sum_{r=0}^\infty\sum_{\substack{x\in Q_1\\d(x,\partial Q_1)=r}}\sum_{\substack{\Gamma\ni\tilde x\\\diam(\Gamma)>r}}\|\Phi_\Gamma\|_\infty\,.
    \end{equation}
    By translation invariance of the interaction, the last sum is independent of the choice of $x$, i.e.\ $\sum_{\substack{\Gamma\ni\tilde x\\\diam(\Gamma)>r}}\|\Phi_\Gamma\|_\infty=\sum_{\substack{\Gamma\ni0\\\diam(\Gamma)>r}}\|\Phi_\Gamma\|_\infty$. This allows us to consider the second sum in isolation. It is bounded by
    \begin{equation}
        \sum_{\substack{x\in Q_1\\d(x,\partial Q_i)=r}}\leq|\partial Q_1|\leq Cb^{\alpha n}.
    \end{equation}
    where $\alpha<1$ since the boundary scales strictly less than the volume. Thus, we obtain
    \begin{equation}
        \|\Delta H_n^\Phi\|_\infty\leq b\tilde C\cdot b^{\alpha n}\sum_{r=0}^\infty\sum_{\substack{\Gamma\ni0\\\diam(\Gamma)>r}}\|\Phi_\Gamma\|_\infty\,.
    \end{equation}
    For the remaining sum, we note that each $\Gamma$ appears exactly $\diam(\Gamma)$ times such that
    \begin{equation}
        \sum_{r=0}^\infty\sum_{\substack{\Gamma\ni0\\\diam(\Gamma)>r}}\|\Phi_\Gamma\|_\infty=\sum_{\Gamma\ni0}\diam(\Gamma)\cdot\|\Phi_\Gamma\|_\infty=M<\infty.
    \end{equation}
    Hence, we conclude
    \begin{equation}
        \|\Delta H_n^\Phi\|_\infty\leq bM\tilde C\cdot b^{\alpha n}\implies H^\Phi\in \mathcal{V}_\phys.
    \end{equation}
\end{proof}

\end{document}
